\chapter*{Список использованных источников}

\begin{enumerate}
	\item Ануфриев Б.Ф. Курс лекций по дисциплине "Техника и методы физических измерений и методы физических расчетов". НИЯУ МИФИ, 2025.
	
	\item Таганцев, А. К. Пиро-, пьезо-, флексоэлектрический и термополяризационный эффекты в ионных кристаллах // Успехи физических наук. 1987. Т. 152, № 3. С. 537.
	
	\item Adams, Thomas M. and Layton, Richard A. Introductory MEMS. Fabrication and Applications. New York, NY: Springer, 2010. ISBN 978-0-387-09510-3. DOI: 10.1007/978-0-387-09511-0.

	\item Бабаева Н.Ф. Гироскопы. Л.: Машиностроение, 1973. 104 с.
	
	\item Шереметьев А.Г. Волоконно-оптический гироскоп. М.: Радио и связь, 1987. 152 с.
	
	\item Малыкин Г.Б. Эффект Саньяка. Корректные и некорректные объяснения // УФН. 2000. Т. 172, № 6.

	\item Филатов Ю.В. Оптические гироскопы. СПб.: Государственный научный центр Российской Федерации ЦНИИ "Электроприбор", 2005.
	
	\item Кашкаров А.П. Микроэлектромеханические системы и элементы. М.: ДМК Пресс, 2018. 114 с.

	\item АО "Фирма ТВЭЛА" [Электронный ресурс]. URL: \url{https://www.tvema.ru/645} (дата обращения: 20.01.2026).
	
	\item АО НЦП ИНФОРТАНС [Электронный ресурс]. URL: \url{https://infotrans-logistic.ru/tiwir03} (дата обращения: 20.01.2026).
	
	\item Описание системы измерения износа рельсов / ЗАО «ПИК Прогресс». 2017. – Лазерный измерительный комплекс для контроля профиля рельса. Состав: датчики износа рельса (ИР), датчик ширины колеи (ШК), вычислительный комплекс ЦНИИ-4МД, датчик пути.
\end{enumerate}